\documentclass[11pt]{article}

\usepackage[T1]{fontenc}
\usepackage[utf8]{inputenc}
\usepackage[english]{babel}
\usepackage[margin=1in]{geometry}
\usepackage{lmodern}
\usepackage{microtype}
\usepackage{mathtools}
\usepackage{amsmath,amssymb,amsthm}
\usepackage{array}
\usepackage{booktabs}
\usepackage{tabularx}
\usepackage{longtable}
\usepackage{float}
\usepackage{graphicx}
\usepackage{pdflscape}
\usepackage{caption}
\usepackage{enumitem}
\usepackage{fancyhdr}
\usepackage[numbers]{natbib}
\usepackage{doi}
\usepackage{xurl}
\usepackage{hyperref}
\usepackage{cleveref}
\usepackage{orcidlink}
\usepackage{titling}



\newcommand{\PaperOne}{Paper~1}
\newcommand{\PaperTwo}{Paper~2}

\newcommand{\CSL}{\textsc{csl}}            \newcommand{\CSLcore}{Candidate Structural Law}    

\newcommand{\Substrate}{\sigma}             \newcommand{\Lift}{L}                       \newcommand{\LiftState}{\sigma_{L}}         \newcommand{\HistorySpace}{H_{\Substrate}^{\Lift}}

\newcommand{\Qquotient}{Q_{\Substrate}^{\Lift}}
\newcommand{\Mquotient}{M_{\Substrate}^{\Lift}}
\newcommand{\ForgetWitnessMap}{\pi}
\newcommand{\Qclass}[1]{[#1]_{Q}}
\newcommand{\Mclass}[1]{[#1]_{M}}

\newcommand{\CurrentObservable}{\varphi_{Q}}    \newcommand{\Exposure}{\operatorname{Exposure}} 

\newcommand{\TraceFPLoopClass}{L_{\mathrm{traceFP}}}
\newcommand{\LoopActionQ}[2]{#1\cdot\Qclass{#2}}
\newcommand{\LoopActionM}[2]{#1\cdot\Mclass{#2}}

\newcommand{\StrengthClass}[1]{S_{#1}}      

\newcommand{\PiQP}{\pi_{QP}}                \newcommand{\PiMNP}{\pi_{MNP}}              

\newcommand{\DiagnosisGrade}{\emph{diagnosis grade}}
\newcommand{\RDiagnosis}{\emph{R-Diagnosis}}

\newcommand{\FATCD}{\mathsf{FATCD}}
\newcommand{\BirdInt}{\mathbf{BirdInt}^{\mathrm{aud}}_{\mathrm{fin}}}

\newcommand{\CSLSATHidden}{\Gamma_{\textsc{csl-sat-hidden}}}
 
\newtheorem{theorem}{Theorem}[section]
\newtheorem{lemma}[theorem]{Lemma}
\newtheorem{proposition}[theorem]{Proposition}
\newtheorem{corollary}[theorem]{Corollary}
\theoremstyle{definition}
\newtheorem{definition}[theorem]{Definition}
\newtheorem{remark}[theorem]{Remark}
\floatstyle{ruled}
\newfloat{algorithm}{tbp}{loa}
\floatname{algorithm}{Algorithm}

\urlstyle{same}
\captionsetup{
  font=small,
  labelfont=bf,
  labelsep=period
}
\captionsetup[table]{skip=6pt,position=top}
\captionsetup[figure]{skip=6pt}
\captionsetup[algorithm]{skip=6pt,position=top}
\crefname{algorithm}{algorithm}{algorithms}
\Crefname{algorithm}{Algorithm}{Algorithms}
\crefname{theorem}{theorem}{theorems}
\Crefname{theorem}{Theorem}{Theorems}
\crefname{lemma}{lemma}{lemmas}
\Crefname{lemma}{Lemma}{Lemmas}
\crefname{proposition}{proposition}{propositions}
\Crefname{proposition}{Proposition}{Propositions}
\crefname{corollary}{corollary}{corollaries}
\Crefname{corollary}{Corollary}{Corollaries}
\crefname{definition}{definition}{definitions}
\Crefname{definition}{Definition}{Definitions}
\crefname{remark}{remark}{remarks}
\Crefname{remark}{Remark}{Remarks}
\setlist[enumerate]{leftmargin=*, itemsep=2pt, topsep=4pt}
\setlength{\droptitle}{-3.2em}
\pretitle{\begin{center}\LARGE\bfseries}
\posttitle{\par\end{center}\vspace{0.5em}}
\preauthor{\begin{center}\normalsize}
\postauthor{\par\end{center}\vspace{-0.4em}}
\predate{\begin{center}\small}
\postdate{\par\end{center}\vspace{0.6em}}
\setlength{\emergencystretch}{2em}
\clubpenalty=10000
\widowpenalty=10000
\displaywidowpenalty=10000
\interfootnotelinepenalty=10000
\allowdisplaybreaks[2]
\fancypagestyle{firstpage}{\fancyhf{}
  \fancyfoot[C]{\scriptsize Automorph Inc., Wilmington, DE, USA \quad \texttt{ioannis@automorph.io} \quad \textcopyright\ 2026 \quad CC-BY 4.0}
  \renewcommand{\headrulewidth}{0pt}
  \renewcommand{\footrulewidth}{0pt}
}

\title{Hiddenness as a Structural Law of Emergence}
\author{Ioannis Tsiokos\,\orcidlink{0009-0009-7659-5964}}
\date{\today}
\hypersetup{
  pdftitle={Hiddenness as a Structural Law of Emergence},
  pdfauthor={Ioannis Tsiokos},
  hidelinks
}

\begin{document}

\maketitle
\thispagestyle{firstpage}

\begin{abstract}
When does prediction become trivial?  This paper gives a typed
structural answer: predictive content collapses exactly when an
upstream determining state is exposed to current observation.  The
load-bearing variable is the visibility coupling predicate
\(\Exposure(\Lift,\Substrate)\), which asks whether some polynomial-time
current observable factors through the autonomous lift's internal state
by a function non-constant on the realized carrier.  We state a typed
\CSLcore{} (\(\CSL\)) over the predictive-quotient holonomy carrier on
autonomous-lifted polynomial-time substrates.  The law has four clauses:
dissolution under exposure, preservation under non-exposure, a
degenerate clause when there is no predictive witness content, and
strength-independence of the coupling function.  The \CSL{} is stated
at \DiagnosisGrade{}.  We also prove a granularity-mismatch theorem at
\RDiagnosis{}: the framework's current and predictive quotients are
strictly finer than the standard classes \(P\) and \(NP\) at the
history level, but this strictness can be lost under standard
language-level projection.
\end{abstract}

\medskip
\noindent\textbf{Keywords.} Six Birds Theory; predictive-quotient holonomy; visibility coupling; exposure predicate; Candidate Structural Law; granularity mismatch; computational complexity; structural law.



\section{Introduction}
\label{sec:intro}

When does prediction become trivial?  In an ordinary computation one
can distinguish the information visible at the present time from the
information that would become visible only after consulting a witness,
a continuation, or a later trace.  Sometimes this distinction is real:
two computation histories look identical to every current test, yet
they separate once witness-augmented tests are allowed.  Sometimes the
distinction is only apparent: the purportedly predictive information is
already encoded in a state variable that a current polynomial-time
observer can read.  The question of this paper is where that boundary
lies.

The answer proposed here is structural.  Prediction
collapses exactly when the upstream state that determines the future
behavior is exposed to current observation.  The strength of a feedback
function, seed schedule, or coupling map is not the primary variable:
a very strong hidden coupling can preserve predictive content, while a
weak but visible coupling can dissolve it.  What matters is whether a
polynomial-time current observable can factor through the internal
determining state in a way that is non-constant on the states actually
realized by the bounded carrier.  This condition is the visibility
coupling predicate, written \(\Exposure(\Lift,\Substrate)\).

The language used to make this precise comes from Six Birds Theory, a
typed, audited calculus of emergence.  The present paper applies it to
a computational substrate, where we use it as a vocabulary for naming
structural data: what the carrier is, which observables are allowed,
what is recorded by an audit, and how quotients are formed.  The underlying finite-audit machinery is the
\(\FATCD\) discipline of Foundations II and the finite audited
interaction calculus \(\BirdInt\) of Foundations III; in ordinary terms,
these provide typed records saying what data are present, what
projections are lawful, and what visibility tags have been declared
\citep{TsiokosFoundationsII2026,TsiokosFoundationsIII2026}.

The basic carrier is a \emph{substrate} \(\Substrate\): a typed bounded
context specifying an instance family, a size bound, and the trace
structure by which computation histories are generated.  An
\emph{autonomous lift} \(\Lift\) is a deterministic policy over that
substrate together with an internal-state readout
\(\LiftState\).  In standard language, it is a self-contained
generative process whose behavior is determined by its state.  For
such a pair \((\Substrate,\Lift)\), the history space
\(\HistorySpace\) is the finite typed set of bounded polynomial-time
histories generated by \(\Lift\) on \(\Substrate\).  A history records
the input instance, the solver-prefix trace, the lift state
\(\LiftState(h)\), declared continuations, and audit bookkeeping.
This is the predictive-quotient holonomy carrier developed in the
earlier holonomy-with-memory work \citep{TsiokosHolonomyMemory2026}.

On this same carrier there are two quotients.  The \emph{current
quotient} \(\Qquotient\) identifies two histories when every declared
polynomial-time current observable gives the same value on them.  The
\emph{predictive quotient} \(\Mquotient\) identifies two histories
when every declared witness-augmented polynomial-time observable gives
the same value under the admissible continuations.  Thus
\(\Mquotient\) remembers witness-side and continuation-side structure
that \(\Qquotient\) may forget.  Predictive witness content is the
existence of two histories that are equal in \(\Qquotient\) but
distinct in \(\Mquotient\): the carrier then contains predictive
information beyond what current observation exposes.

The central object of the paper is the \CSLcore{} (\(\CSL\)).  It is a
four-clause structural law over the carrier just described.  In plain
terms, the law says: if exposure holds and predictive witness content
is present, then the apparent extra predictive motion dissolves into
current motion; if exposure fails and predictive witness content is
present, then the M-versus-Q asymmetry is preserved; if there is no
predictive witness content, there is no asymmetry to explain; and the
dissolution status depends only on the binary exposure predicate, not
on a separate strength scale for the coupling function.  The formal
statement in \Cref{sec:csl-core} gives the typed hypotheses under
which this four-clause shape is proved.

The second headline result records a boundary with standard
complexity-theoretic language.  The quotients \(\Qquotient\) and
\(\Mquotient\) are quotients of histories: they remember trace
features, witness identity, solver-prefix behavior, and audit
bookkeeping.  The standard classes \(P\) and \(NP\), by contrast, are
language classes.  The natural projections from \(\Qquotient\) to
polynomial-time decision predicates and from \(\Mquotient\) to
existential verifier predicates discard trace and witness identity.
The granularity-mismatch theorem in
\Cref{sec:granularity-mismatch} exhibits this loss explicitly: strict
refinement at the history-quotient level can collapse after projection
to language-level predicates.

The scope of the paper is correspondingly narrow.  The \CSL{} is a
framework-internal structural law at \DiagnosisGrade{}, meaning
bounded finite-carrier evidence with a falsifier audit, not an
asymptotic theorem over all polynomial-time computation histories.  The
granularity-mismatch theorem is at \RDiagnosis{}, meaning that it
identifies a structural obstruction while the representation theorem
from arbitrary polynomial-time deciders to these autonomous-lift
carriers is named but not constructed.  This paper does not claim that
the \CSL{} implies a standard \(P\neq NP\) separation
\citep{Cook2000PNPClay,Aaronson2016PNP}; the precise scope is recorded
in \Cref{sec:scope}.
 

\section{Framework}
\label{sec:framework}

We now make precise the typed objects introduced in
\Cref{sec:intro}.  This section first fixes the computational
substrate and the finite history carrier; it then places two
equivalence relations on that carrier --- the current and predictive
quotients --- and records the refinement map between them.

The bookkeeping is typed in the sense of the underlying finite-audit
framework: the instance family, size bound, trace data, state record,
and audit record are all declared as part of the object rather than
left implicit.  This is the role played by the \(\FATCD\) discipline
and the finite audited interaction calculus \(\BirdInt\) in the
background foundations \citep{TsiokosFoundationsII2026,TsiokosFoundationsIII2026}.

\begin{definition}[Computational substrate]
\label{def:substrate}
A \emph{computational substrate} \(\Substrate\) is a typed bounded
polynomial-time computation context \citep{AroraBarak2009Complexity}.
It consists of:
\begin{enumerate}
  \item an instance family \(X_{\Substrate}\);
  \item a size function on instances together with a bound
  \(N_{\Substrate}\), so that only instances of size at most
  \(N_{\Substrate}\) belong to the bounded carrier considered here;
  \item a declared solver-prefix family, assigning to each bounded
  instance the finite trace shape that a polynomial-time computation
  may generate before the history is audited.
\end{enumerate}
A substrate is therefore the typed context in which instances, scale,
and trace formation have all been fixed, rather than a language or a
set of inputs.
\end{definition}

\begin{remark}[Example: 3-SAT substrate]
Fix a bound \(N\).  One may take \(X_{\Substrate}\) to be the family
of 3-CNF formulas of size at most \(N\), measured for instance by the
number of variables and clauses
\citep{Cook1971TheoremProving,Karp1972Reducibility,GareyJohnson1979Intractability}.  The solver-prefix family can be a
declared polynomial-time decision-attempt policy: for each formula it
records the finite sequence of inspected clauses, partial assignments,
branching choices, local simplifications, and time-certificate data
produced before the bounded audit closes.  No completeness claim for
3-SAT is built into this example; the substrate only declares the
finite typed traces on which the later quotients will act.
\end{remark}

The next object is the carrier on which current and predictive
observations will be compared.  An autonomous lift \(\Lift\) on
\(\Substrate\) is, for the purposes of this section, a deterministic
polynomial-time policy together with an internal-state readout
\(\LiftState\); its full definition is given in
\Cref{sec:autonomous-lifts}.

\begin{definition}[History space]
\label{def:history-space}
Let \(\Substrate\) be a computational substrate and let \(\Lift\) be
an autonomous lift on \(\Substrate\).  The \emph{history space}
\(\HistorySpace\) is the finite typed set of bounded polynomial-time
computation histories generated by \(\Lift\) on \(\Substrate\).  Each
history \(h\in\HistorySpace\) records:
\begin{enumerate}
  \item the substrate input instance;
  \item the solver-prefix trace generated under \(\Lift\);
  \item the autonomous-lift internal-state record \(\LiftState(h)\);
  \item the declared predictive continuations available from \(h\);
  \item the audit bookkeeping attached to \(h\).
\end{enumerate}
\end{definition}

This finite carrier is the bounded form of the predictive-quotient
holonomy carrier used in the earlier holonomy-with-memory framework
\citep{TsiokosHolonomyMemory2026}.  The next step is to place two
equivalence relations on \(\HistorySpace\)
\citep{Halmos1960NaiveSetTheory}: one for what current polynomial-time
observation can see, and one for what witness-augmented prediction can
still distinguish.

A \emph{current observable} is a declared polynomial-time function on
\(\HistorySpace\) whose value is available at the current stage of the
history: decision output, visible trace features, solver-prefix
bookkeeping, or other audited quantities not requiring a witness or a
future continuation.  We write \(\CurrentObservable\) for such an
observable when its codomain is understood.

\begin{definition}[Current quotient]
\label{def:current-quotient}
The \emph{current quotient} \(\Qquotient\) is the equivalence relation
on \(\HistorySpace\) that identifies two histories exactly when every
declared polynomial-time current observable gives the same value on
them.  Thus, for \(h,h'\in\HistorySpace\),
\[
  \Qclass{h}=\Qclass{h'}
  \quad\Longleftrightarrow\quad
  \CurrentObservable(h)=\CurrentObservable(h')
  \text{ for every declared current observable }\CurrentObservable .
\]
\end{definition}

A \emph{witness-augmented observable} is a declared polynomial-time
test allowed to inspect the witness-side data and the admissible
continuation records attached to a history.  In ordinary terms, it is
the same kind of finite typed test as a current observable, but with
the additional witness and continuation fields left visible.

\begin{definition}[Predictive quotient]
\label{def:predictive-quotient}
The \emph{predictive quotient} \(\Mquotient\) is the equivalence
relation on \(\HistorySpace\) that identifies two histories exactly
when every declared witness-augmented polynomial-time observable agrees
on them under all admissible continuations.  Thus, for
\(h,h'\in\HistorySpace\),
\[
  \Mclass{h}=\Mclass{h'}
  \quad\Longleftrightarrow\quad
  \varphi_{M}(h)=\varphi_{M}(h')
  \text{ for every declared witness-augmented observable }\varphi_{M}.
\]
\end{definition}

The predictive quotient is therefore the finer comparison whenever the
witness-side fields carry information not visible to current
observation.  The next definition records the canonical way to forget
that extra witness-side information.

\begin{definition}[Forget-witness map]
\label{def:forget-witness-map}
The \emph{forget-witness map}
\[
  \ForgetWitnessMap:\Mquotient\longrightarrow\Qquotient
\]
sends each predictive class to the unique current class containing it:
\[
  \ForgetWitnessMap(\Mclass{h})=\Qclass{h}.
\]
It is well defined because every polynomial-time current observable is
also a witness-augmented observable with the witness and continuation
fields ignored.
\end{definition}

The carrier has predictive witness content, in the sense formalized in
\Cref{sec:witness-content}, when \(\Mquotient\) strictly refines
\(\Qquotient\) on the support of \(\HistorySpace\).

\begin{proposition}[Witness refinement]
\label{prop:witness-refinement}
The predictive quotient refines the current quotient: for all
\(h,h'\in\HistorySpace\), if \(\Mclass{h}=\Mclass{h'}\), then
\(\Qclass{h}=\Qclass{h'}\).  Equivalently, the map
\(\ForgetWitnessMap:\Mquotient\to\Qquotient\) is surjective.
\end{proposition}

\begin{proof}
Unfold the two equivalence relations.  Suppose
\(\Mclass{h}=\Mclass{h'}\).  Then every declared
witness-augmented polynomial-time observable gives the same value on
\(h\) and \(h'\), under all admissible continuations.  A current
observable \(\CurrentObservable\) is a special case: view it as a
witness-augmented observable that ignores the witness and continuation
coordinates.  Hence \(\CurrentObservable(h)=\CurrentObservable(h')\)
for every declared current observable \(\CurrentObservable\), which is
exactly \(\Qclass{h}=\Qclass{h'}\).  Surjectivity follows because
every current class \(\Qclass{h}\) is the image of the predictive class
\(\Mclass{h}\) under \(\ForgetWitnessMap\).\footnote{This proposition
is verified in Lean by the witness-refinement declaration listed in
\Cref{app:formalization}.}
\end{proof}
 

\section{Autonomous lifts}
\label{sec:autonomous-lifts}

The carrier \(\HistorySpace\) depends not only on the substrate
\(\Substrate\), but also on the rule by which histories are generated.
We now define that rule precisely as an autonomous lift, and then
record the strength stratification used later in the
strength-independence clause of the \CSL{}.

\begin{definition}[Autonomous lift]
\label{def:autonomous-lift}
Let \(\Substrate\) be a computational substrate.  An
\emph{autonomous lift} on \(\Substrate\) is a typed object
\[
  \Lift=(\mathrm{gen}_{\Lift},\LiftState,\mathrm{replay}_{\Lift})
\]
with the following components:
\begin{enumerate}
  \item \(\mathrm{gen}_{\Lift}\) is a deterministic polynomial-time
  history-generation procedure.  In plain terms, it is a deterministic
  polynomial-time algorithm that, given the substrate input and the
  prior trace data, extends the computation history by one audited
  step.
  \item \(\LiftState:\HistorySpace\to S\) is the internal-state
  readout of the lift, where \(S\) is the lift's finite typed state
  space.  The value \(\LiftState(h)\) records the internal determining
  state of the procedure at the history \(h\).
  \item \(\mathrm{replay}_{\Lift}\) is the all-input replay record for
  \(\mathrm{gen}_{\Lift}\): an audited machine-index and
  time-certificate object that allows the same polynomial-time
  procedure to be reproduced on any input.  This is the audit-side
  bookkeeping supplied by the \(\FATCD\) discipline
  \citep{TsiokosFoundationsII2026}.\footnote{The autonomous lift is
  realized in Lean by the autonomous-lift declaration listed in
  \Cref{app:formalization}; the Lean record exposes additional
  family-characterizing predicates beyond the components named in the
  definition above.  The narrowing-disclosure discipline is recorded in
  \Cref{app:formalization}.}
\end{enumerate}
\end{definition}

In standard mathematical language, an autonomous lift is a deterministic
finite-control algorithm whose behavior depends only on its declared
inputs, its prior trace, and its own internal state.  The readout
\(\LiftState\) exposes that internal state as a function on histories.
The replay record is what distinguishes the audited lift from an
unaudited algorithm with the same input-output behavior: it records
which machine is being run and why the bounded run is accepted as a
lawful polynomial-time history generator.

\begin{remark}[Strength stratification]
The later \CSL{} statement refers to a four-level strength
stratification
\[
  \StrengthClass{1},\ \StrengthClass{2},\ \StrengthClass{3},\
  \StrengthClass{4}
\]
for coupling functions used inside \(\mathrm{gen}_{\Lift}\).  The
levels should be read as a graded family of implementation strengths:
\(\StrengthClass{1}\) is trivial or absent coupling,
\(\StrengthClass{2}\) is weak coupling such as linear feedback,
\(\StrengthClass{3}\) is standard cryptographic-strength coupling in
the sense of pseudorandom-function families
\citep{GoldreichGoldwasserMicali1986RandomFunctions,Goldreich2001FoundationsCryptography},
and \(\StrengthClass{4}\) is beyond-cryptographic-strength coupling
\citep{HastadImpagliazzoLevinLuby1999PRG}.
These labels are not formalized here as complexity classes.  They are
placeholders for any fixed four-level grading of coupling strength.
The strength-independence clause in \Cref{sec:csl-core} says that the
dissolution status of the M-versus-Q asymmetry depends only on the
binary exposure predicate, not on which class
\(\StrengthClass{k}\) contains the coupling function.
\end{remark}

Four canonical families of autonomous lifts illustrate the structural
distinctions that the \CSL{} uses.  They name recurring ways in which
\(\mathrm{gen}_{\Lift}\) can depend on state, trace, randomness, and
witness-side information: a self-generating family, a feedback-coupled
family, an internally seeded random family, and a planted-suppressed
family.

\begin{definition}[Lift A---Self-generating substrate]
\label{def:lift-family-A}
An autonomous lift \(\Lift\) on \(\Substrate\) is in \emph{Lift A} if
\(\mathrm{gen}_{\Lift}\) generates substrate instances internally from
the lift's own state, without an external instance feed.  In plain
terms, the lift is itself the instance generator: for example, a
deterministic algorithm that, given only its current internal state,
emits a new bounded 3-SAT instance \citep{Cook1971TheoremProving} and
then attempts to decide it.
\end{definition}

\begin{definition}[Lift B---Feedback-coupled]
\label{def:lift-family-B}
An autonomous lift \(\Lift\) on \(\Substrate\) is in \emph{Lift B} if
\(\mathrm{gen}_{\Lift}\) couples instance generation to the running
trace, so that the lift's internal state and prior trace data feed
back into the next instance-generation step.  In plain terms, the
lift produces the next instance as a function of what it has already
done; the strength of that function may range over the
\(\StrengthClass{k}\) stratification.
\end{definition}

\begin{definition}[Lift C---Internally-seeded random]
\label{def:lift-family-C}
An autonomous lift \(\Lift\) on \(\Substrate\) is in \emph{Lift C} if
\(\mathrm{gen}_{\Lift}\) uses an internal seed to inject randomness
into instance generation, while the seed is not visible to current
observation.  Equivalently, the apparent randomness is controlled by a
state variable seen by \(\LiftState\), not by an external oracle; the
current observation channel cannot read that seed directly.
\end{definition}

\begin{definition}[Lift D---Planted-suppressed internalization]
\label{def:lift-family-D}
An autonomous lift \(\Lift\) on \(\Substrate\) is in \emph{Lift D} if
\(\mathrm{gen}_{\Lift}\) plants witness content into the generated
instance but suppresses its visibility from current observation.  For
example, the lift may construct a 3-SAT instance from a satisfying
assignment that it knows internally, while the current observation
channel is configured so that the planted assignment is not exposed.
\end{definition}

These families are not mutually exclusive.  A single autonomous lift
may be both internally seeded and feedback-coupled, or may combine
planted witness content with self-generation.  The names mark the
structural distinction relevant to the exposure predicate: in the
typical configurations, Lift A and Lift D keep the determining state
or witness-side content unavailable to current observation, while Lift
B and Lift C can make that state available depending on the coupling
function and seed exposure.  The \CSL{} states precisely when such
access dissolves the M-versus-Q asymmetry; that statement appears in
\Cref{sec:csl-core}.
 

\section{Loop class and loop actions}
\label{sec:loop-action}

The \CSL{} is formulated in terms of how a designated family of
structural loops interacts with the two quotients on \(\HistorySpace\).
We now fix that loop class, define its actions on \(\Qquotient\) and
\(\Mquotient\), and name the resulting M-versus-Q asymmetry predicate.

Here a \emph{loop class} should be read in concrete finite-carrier
terms.  It is a designated collection of structural rewrites on
histories: relabelings, controlled trace permutations, audited
continuations, or similar transformations that leave the typed carrier
intact \citep{DummitFoote2004AbstractAlgebra}.  The word ``trace-fingerprint'' records that the rewrite may
alter the audited record of trace structure while preserving the
underlying substrate instance and the compatibility data needed to
descend to the quotients.

\begin{definition}[Trace-fingerprint loop class]
\label{def:trace-fingerprint-loop-class}
A \emph{trace-fingerprint loop class}
\(\TraceFPLoopClass\) on \(\HistorySpace\) is a declared family of
loops in the predictive-quotient holonomy carrier.  Each
\(\ell\in\TraceFPLoopClass\) acts on histories by a structured
trace-level modification
\[
  h\longmapsto \ell.h .
\]
The map \(h\mapsto\ell.h\) preserves the finite carrier's typed
structure, preserves the underlying substrate instance, and is
measure-preserving when \(\HistorySpace\) is read with its finite
counting support.  The family is chosen so that its actions are
compatible with both \(\Qquotient\) and \(\Mquotient\).
\end{definition}

\begin{definition}[Loop action on \(Q\)]
\label{def:loop-action-q}
Let \(\ell\in\TraceFPLoopClass\) and let \(\Qclass{h}\in\Qquotient\).
The \emph{loop action on \(Q\)} is
\[
  \LoopActionQ{\ell}{h}=\Qclass{\ell.h}.
\]
This is well defined because the loop class is compatible with
Q-equivalence: if \(\Qclass{h}=\Qclass{h'}\), then
\(\Qclass{\ell.h}=\Qclass{\ell.h'}\).
\end{definition}

\begin{definition}[Loop action on \(M\)]
\label{def:loop-action-m}
Let \(\ell\in\TraceFPLoopClass\) and let \(\Mclass{h}\in\Mquotient\).
The \emph{loop action on \(M\)} is
\[
  \LoopActionM{\ell}{h}=\Mclass{\ell.h}.
\]
This is well defined because the loop class is compatible with
M-equivalence: if \(\Mclass{h}=\Mclass{h'}\), then
\(\Mclass{\ell.h}=\Mclass{\ell.h'}\).
\end{definition}

\begin{definition}[M-vs-Q loop-action asymmetry]
\label{def:asymmetry}
The pair \((\Substrate,\Lift)\) exhibits an
\emph{M-vs-Q loop-action asymmetry} if there exist
\(\ell\in\TraceFPLoopClass\) and \(h\in\HistorySpace\) such that the
loop action is trivial on the current class but non-trivial on the
predictive class:
\[
  \LoopActionQ{\ell}{h}=\Qclass{h}
  \qquad\text{and}\qquad
  \LoopActionM{\ell}{h}\neq\Mclass{h}.
\]
Equivalently, there is a structural rewrite that current observation
cannot detect, because it leaves the \(Q\)-class fixed, while
predictive observation can detect it, because it moves the
\(M\)-class.
\end{definition}

The asymmetry predicate is a yes-or-no condition attached to the
chosen pair \((\Substrate,\Lift)\) and its declared loop class.  When
it holds, the M-versus-Q distinction is preserved by a loop that is
invisible to current observation.  When predictive witness content is
present and this asymmetry fails, the paper calls that failure
dissolution: the predictive motion has no independent quotient-level
effect beyond what current observation already captures.  The next
sections define the exposure predicate and then state the CSL relation
between exposure, preservation, and dissolution.
 

\section{Visibility coupling}
\label{sec:visibility-coupling}

We now define the visibility coupling predicate
\(\Exposure(\Lift,\Substrate)\).  This is the load-bearing predicate
of the \CSL{}: the four-clause structural law in
\Cref{sec:csl-core} is organized entirely around whether exposure
holds or fails.

The predicate is meant to capture a simple standard mathematical
condition.  The lift has an internal-state readout
\(\LiftState:\HistorySpace\to S\), where \(S\) is the lift's finite
state space.  Exposure holds when some polynomial-time current
observable can read a non-trivial function of that internal state.  In
other words, current observation does not merely see the instance or
the visible trace; it detects a genuine variation in the lift's
determining state.

Two nearby conditions are excluded.  First, the lift state
may vary while no declared polynomial-time current observable can read
that variation; then there is no exposure.  Second, a function
\(f:S\to\mathcal O\) may be non-constant on the abstract state space
\(S\) but constant on the states that actually occur in histories; it
also does not witness exposure.  The predicate is about what is visible
on the realized carrier, not about irrelevant state values.

\begin{definition}[Visibility coupling / exposure]
\label{def:visibility-coupling}
Let \(\Substrate\) be a computational substrate, let \(\Lift\) be an
autonomous lift on \(\Substrate\), and write
\(\LiftState:\HistorySpace\to S\) for its internal-state readout.
The \emph{visibility coupling predicate}
\(\Exposure(\Lift,\Substrate)\) holds if and only if there exist an
observable codomain \(\mathcal O\), a declared polynomial-time current
observable
\[
  \CurrentObservable:\HistorySpace\to\mathcal O,
\]
and a polynomial-time function \(f:S\to\mathcal O\) such that
\[
  \bigl(\exists h,h'\in\HistorySpace,\,
      f(\LiftState(h))\neq f(\LiftState(h'))\bigr)
  \quad\text{and}\quad
  \bigl(\forall h\in\HistorySpace,\,
      \CurrentObservable(h)=f(\LiftState(h))\bigr).
\]
\end{definition}

The first conjunct is the realized-support nonconstancy condition.  It
says that \(f\) is non-constant on the realized image
\[
  \{\LiftState(h)\mid h\in\HistorySpace\},
\]
not merely on the ambient state space \(S\).  The second conjunct says
that this non-trivial state function is actually visible through a
polynomial-time current observable on the history carrier.

\begin{remark}[Example: planted witness without exposure]
Consider a Lift D-like configuration in which the internal state space
contains a planted witness, say \(S=\{0,1\}^{k}\), and
\(\LiftState(h)\) records the planted assignment used to generate a
bounded 3-SAT instance.  The abstract function \(f:S\to S\) given by
the identity is certainly non-constant, and it may even be
non-constant on the realized states.  Nevertheless, if the declared
current observables see only the formula, visible trace features, and
audit bookkeeping, and no polynomial-time current observable computes
that planted assignment on \(\HistorySpace\), then
\(\Exposure(\Lift,\Substrate)\) fails.  The state carries predictive
content, but the current observation channel is configured to suppress
it.
\end{remark}

\begin{remark}[Example: feedback exposure]
Consider a Lift B-like configuration in which a feedback-coupled
generator stores an internal seed coordinate in \(\LiftState(h)\), and
that coordinate takes at least two values on realized histories.  If
the decision-output observable depends non-trivially on that seed
coordinate, then exposure holds.  Take \(\CurrentObservable\) to be
the decision-output observable and take \(f\) to be the seed-coordinate
projection composed with the decision rule on that coordinate.  The
factorization \(\CurrentObservable(h)=f(\LiftState(h))\) holds on
\(\HistorySpace\), and the seed dependence gives nonconstancy on the
realized image.
\end{remark}

The contrast is the point of the definition.  Lift D-like
configurations can carry rich witness-side information while satisfying
\(\neg\Exposure(\Lift,\Substrate)\), because the visibility channel
suppresses the determining state.  Lift B-like configurations often
satisfy \(\Exposure(\Lift,\Substrate)\), because their feedback
structure produces an observable that factors through the lift state.
The \CSL{} will say that, when predictive witness content is present,
exposure dissolves the M-versus-Q asymmetry, while non-exposure
preserves it.

Finally, \(\Exposure(\Lift,\Substrate)\) is finite-carrier testable.
It ranges over the finite history space \(\HistorySpace\), the
realized image of \(\LiftState\) on that carrier, and the declared
families of polynomial-time current observables and polynomial-time
state functions.  No mutual-information estimate or algebraic
dependence test is part of the predicate; the functional-dependence
condition is direct.
 

\section{Predictive witness content}
\label{sec:witness-content}

The \CSL{} distinguishes dissolution from preservation only when there
is something for that distinction to act on.  We now name this
precondition: the selected carrier has predictive witness content when
the predictive quotient strictly refines the current quotient on the
support of \(\HistorySpace\).

\begin{definition}[Predictive witness content]
\label{def:predictive-witness-content}
The pair \((\Substrate,\Lift)\) has \emph{predictive witness content}
in the selected carrier if there exist histories
\(h,h'\in\HistorySpace\) such that
\[
  \Qclass{h}=\Qclass{h'}
  \qquad\text{but}\qquad
  \Mclass{h}\neq\Mclass{h'} .
\]
Equivalently, the forget-witness map
\[
  \ForgetWitnessMap:\Mquotient\to\Qquotient
\]
is not a bijection: \(\Mquotient\) strictly refines \(\Qquotient\) on
the support of \(\HistorySpace\).
\end{definition}

In standard mathematical terms, predictive witness content means that
two histories look identical to every current observable but separate
once witness-augmented observation is allowed.  The presence of such a
pair is the operational meaning of saying that the predictive quotient
contains more information than the current quotient.

\begin{remark}[Example: two satisfying witnesses]
A prototypical case is a bounded 3-SAT instance
\citep{Cook1971TheoremProving,GareyJohnson1979Intractability} with two
distinct satisfying assignments.  Two histories may have the same visible
instance and the same decision-bit current output, and therefore agree
under decision-level current observation.  A witness-augmented
observable that reads the satisfying assignment can still distinguish
the histories.  In the notation above, this gives
\(\Qclass{h}=\Qclass{h'}\) but
\(\Mclass{h}\neq\Mclass{h'}\).
\end{remark}

If \((\Substrate,\Lift)\) lacks predictive witness content, then
\(\Mquotient\) and \(\Qquotient\) coincide on the carrier support and
\(\ForgetWitnessMap\) is a bijection.  In that case the
dissolution/preservation distinction is vacuous: there is no
M-versus-Q asymmetry to preserve or dissolve.  The degenerate clause
of the \CSL{} records exactly this case, and the next section states
the full structural law.
 

\section{CSL core lemma}
\label{sec:csl-core}

We now state the \CSLcore{} (\(\CSL\)).  The law has four clauses:
dissolution, preservation, the degenerate case, and
strength-independence.  Together they relate the M-vs-Q loop-action
asymmetry of \Cref{sec:loop-action} to the exposure predicate of
\Cref{sec:visibility-coupling} and the predictive-witness precondition
of \Cref{sec:witness-content}.  The law is stated at
\DiagnosisGrade{}, meaning bounded finite-carrier evidence with a
falsifier audit, not an asymptotic claim over all polynomial-time
histories.  The discipline follows the tradition of computer-assisted
finite-case verification used elsewhere in mathematics
\citep{AppelHaken1989FourColorable,HalesEtAl2017Flyspeck,BaileyEtAl2007ExperimentalMath}.

\paragraph{Structural premises required by the typed abstract statement.}
At the typed abstract level, the law assumes the following framework
facts about the predictive-quotient holonomy carrier and the
trace-fingerprint loop class.  These are hypotheses made available to
the abstract statement; they are not consequences of the \CSL{} itself.
The quotient and loop vocabulary is the carrier vocabulary developed
for predictive-quotient holonomy, while the audited state and replay
bookkeeping are typed by the finite-audit framework
\citep{TsiokosHolonomyMemory2026,TsiokosFoundationsII2026}.
\begin{enumerate}[label=(\alph*)]
  \item \emph{Loop-state detection.}  If a loop
  \(\ell\in\TraceFPLoopClass\) changes the predictive class of a
  history, then it also changes the lift state:
  \[
    \LoopActionM{\ell}{h}\neq\Mclass{h}
    \quad\Longrightarrow\quad
    \LiftState(\ell.h)\neq\LiftState(h).
  \]
  This premise records that the loop class is rich enough to move the
  internal state, not merely to relabel quotient classes.

  \item \emph{Exposure-observable detects state change.}  Suppose
  \(\CurrentObservable=f\circ\LiftState\), with \(f\) non-constant on
  the realized image of \(\LiftState\).  If a loop changes
  \(\LiftState(h)\), then it changes the exposed observable value
  \(\CurrentObservable(h)\).  This premise records that the current
  observable actually detects the state variation that exposure claims
  to expose.

  \item \emph{Loop Q-action via observable.}  If a loop changes some
  declared current observable at \(h\), then it changes the current
  quotient class:
  \[
    \CurrentObservable(\ell.h)\neq\CurrentObservable(h)
    \quad\Longrightarrow\quad
    \LoopActionQ{\ell}{h}\neq\Qclass{h}.
  \]
  This clause says that Q-equivalence is exactly the equivalence
  relation seen by all declared current observables.

  \item \emph{Loop preservation witness.}  When
  \(\neg\Exposure(\Lift,\Substrate)\), the loop class contains a
  connector loop for any pair \(h,h'\in\HistorySpace\) with
  \(\Qclass{h}=\Qclass{h'}\) and
  \(\Mclass{h}\neq\Mclass{h'}\).  This premise records that the loop
  class is rich enough to witness predictive separation that current
  observation cannot see.

  \item \emph{M-refines-Q.}  For all histories \(h,h'\in\HistorySpace\),
  \[
    \Mclass{h}=\Mclass{h'}
    \quad\Longrightarrow\quad
    \Qclass{h}=\Qclass{h'} .
  \]
  This is the refinement direction from
  \Cref{prop:witness-refinement}, not a \CSL{} conclusion.

  \item \emph{Strength-decision equivalence.}  The predicate deciding
  the strength-indexed dissolution status is equivalent to
  \(\Exposure(\Lift,\Substrate)\).  This clause says that the strength
  label \(\StrengthClass{k}\) enters through the implementation of
  \(\mathrm{gen}_{\Lift}\), while the law's decision variable is the
  binary exposure predicate.
\end{enumerate}

\begin{lemma}[\CSLcore{} (core)]
\label{lem:csl-core}
Let \(\Substrate\) be a computational substrate, let \(\Lift\) be an
autonomous lift on \(\Substrate\), and let \(\TraceFPLoopClass\) be
the trace-fingerprint loop class on \(\HistorySpace\).  Assume the six
structural premises (a)--(f) above.  Then the following four clauses
hold.
\begin{enumerate}
  \item \emph{Dissolution.}  If \(\Exposure(\Lift,\Substrate)\) holds
  and \((\Substrate,\Lift)\) has predictive witness content, then for
  every loop \(\ell\in\TraceFPLoopClass\) and every history
  \(h\in\HistorySpace\),
  \[
    \Mclass{\ell.h}\neq\Mclass{h}
    \quad\Longrightarrow\quad
    \Qclass{\ell.h}\neq\Qclass{h}.
  \]
  Equivalently, the M-vs-Q loop-action asymmetry fails: every
  M-class change is accompanied by a Q-class change.

  \item \emph{Preservation.}  If \(\Exposure(\Lift,\Substrate)\) fails
  and \((\Substrate,\Lift)\) has predictive witness content, then
  there exist a loop \(\ell\in\TraceFPLoopClass\) and a history
  \(h\in\HistorySpace\) such that
  \[
    \LoopActionQ{\ell}{h}=\Qclass{h}
    \qquad\text{and}\qquad
    \LoopActionM{\ell}{h}\neq\Mclass{h}.
  \]
  The M-vs-Q loop-action asymmetry holds.

  \item \emph{Degenerate.}  If \((\Substrate,\Lift)\) has no
  predictive witness content, then \(\Mquotient=\Qquotient\) on the
  support of \(\HistorySpace\), and the M-vs-Q distinction is vacuous.

  \item \emph{Strength-independence.}  Given the value of
  \(\Exposure(\Lift,\Substrate)\), the dissolution or preservation
  status of the loop-action asymmetry does not depend on the strength
  class \(\StrengthClass{k}\) of the coupling function used inside
  the lift's generation procedure.
\end{enumerate}
\end{lemma}

The lemma is stated at \DiagnosisGrade{}: the empirical hardening of
the law's qualifier --- the falsifier audit over the bounded finite
carrier --- is a paper-side certification rather than part of the
lemma's formal content.

Before the proof, an outline.  The degenerate clause is the negation of
predictive witness content, and strength-independence is a one-line
consequence of exposure being a binary predicate.  The dissolution and
preservation clauses carry the substantive content: they explain how
exposure forces M-motion to become Q-visible, while non-exposure allows
a connector loop to preserve the M-vs-Q asymmetry.

\begin{proof}
\emph{Clause 3, degenerate.}  Suppose first that
\((\Substrate,\Lift)\) has no predictive witness content.  Unfolding
\Cref{def:predictive-witness-content}, this is the negation of the
existence of histories \(h,h'\in\HistorySpace\) with the same current
class and distinct predictive classes.  Equivalently, for all
\(h,h'\in\HistorySpace\),
\[
  \Qclass{h}=\Qclass{h'}
  \quad\Longrightarrow\quad
  \Mclass{h}=\Mclass{h'} .
\]
Premise (e) supplies the refinement direction
\(\Mclass{h}=\Mclass{h'}\Rightarrow\Qclass{h}=\Qclass{h'}\) for the
same histories.  Combining the two implications gives, on the support
of \(\HistorySpace\),
\[
  \Qclass{h}=\Qclass{h'}
  \quad\Longleftrightarrow\quad
  \Mclass{h}=\Mclass{h'}
\]
for every pair \(h,h'\).  The two quotient relations therefore have
the same equivalence classes on the realized carrier.  Hence
\(\Qquotient=\Mquotient\) on that support, and the M-versus-Q
distinction is vacuous.

\emph{Clause 4, strength-independence.}  The predicate
\(\Exposure(\Lift,\Substrate)\) is a binary condition.  It is true
exactly when there exist a declared current observable
\(\CurrentObservable\) and a state function \(f\) satisfying
realized-support nonconstancy and the factorization
\(\CurrentObservable=f\circ\LiftState\) on \(\HistorySpace\); otherwise
it is false.  The strength label \(\StrengthClass{k}\) enters through
the implementation of \(\mathrm{gen}_{\Lift}\), but it is not an
additional value of the exposure predicate.  By premise (f), the
strength-indexed decision predicate for dissolution is equivalent to
\(\Exposure(\Lift,\Substrate)\).  Therefore two implementations with
the same exposure value have the same strength-indexed decision value:
when exposure holds they are in the dissolution branch, and when
exposure fails they are in the preservation branch.  The
dissolution/preservation status is consequently independent of
\(\StrengthClass{k}\) once the value of
\(\Exposure(\Lift,\Substrate)\) is fixed.

\emph{Clause 1, dissolution.}  Assume \(\Exposure(\Lift,\Substrate)\)
holds and that \((\Substrate,\Lift)\) has predictive witness content.
Unfolding predictive witness content gives histories
\(h_{0},h_{1}\in\HistorySpace\) with
\[
  \Qclass{h_{0}}=\Qclass{h_{1}}
  \qquad\text{and}\qquad
  \Mclass{h_{0}}\neq\Mclass{h_{1}} .
\]
Thus the selected carrier is nondegenerate: current observation
identifies histories that witness-augmented observation separates.  An
M-class change along a loop is therefore genuine predictive motion
whose current-quotient visibility is the issue.

By \(\Exposure(\Lift,\Substrate)\), choose a polynomial-time current
observable \(\CurrentObservable:\HistorySpace\to\mathcal O\) and a
polynomial-time function \(f:S\to\mathcal O\) such that \(f\) is
non-constant on the realized image of \(\LiftState\) and
\[
  \CurrentObservable(x)=f(\LiftState(x))
  \qquad\text{for every }x\in\HistorySpace .
\]
Now fix an arbitrary \(\ell\in\TraceFPLoopClass\) and an arbitrary
\(h\in\HistorySpace\), and suppose that
\(\Mclass{\ell.h}\neq\Mclass{h}\).  Equivalently, by the definition of
the loop action on \(M\),
\(\LoopActionM{\ell}{h}\neq\Mclass{h}\).  Premise (a), the loop-state
detection hypothesis, turns this predictive-class change into an
internal-state change:
\[
  \LiftState(\ell.h)\neq\LiftState(h).
\]
Because exposure has supplied a current observable factoring through
\(\LiftState\), premise (b) applies to this state change and gives
\[
  \CurrentObservable(\ell.h)\neq\CurrentObservable(h).
\]
Finally, premise (c) says that any change in a declared current
observable changes the current quotient class.  Therefore
\[
  \LoopActionQ{\ell}{h}=\Qclass{\ell.h}\neq\Qclass{h},
\]
where the equality is the definition of the loop action on \(Q\).
Since \(\ell\) and \(h\) were arbitrary, every M-class change in the
trace-fingerprint loop class is accompanied by a Q-class change.  This
is precisely dissolution of the M-vs-Q loop-action asymmetry.

\emph{Clause 2, preservation.}  Assume
\(\neg\Exposure(\Lift,\Substrate)\) and that \((\Substrate,\Lift)\)
has predictive witness content.  Unfolding the predictive-witness
precondition, choose histories \(h,h'\in\HistorySpace\) such that
\[
  \Qclass{h}=\Qclass{h'}
  \qquad\text{and}\qquad
  \Mclass{h}\neq\Mclass{h'} .
\]
The failure of exposure is exactly the hypothesis needed to invoke
premise (d).  That premise says that, for any pair with this
Q-identification and M-separation, the trace-fingerprint loop class
contains a connector loop.  Hence there is
\(\ell\in\TraceFPLoopClass\) connecting the chosen pair.  Applying the
connector loop to the starting history \(h\), the connector property
keeps the current class fixed while moving the predictive class; using
the definitions of the two loop actions, this gives
\[
  \LoopActionQ{\ell}{h}=\Qclass{h}
  \qquad\text{and}\qquad
  \LoopActionM{\ell}{h}\neq\Mclass{h}.
\]
Thus the loop action is trivial on the current class and non-trivial
on the predictive class, which exhibits the asymmetry-preserving pair.

The four clauses together prove the lemma.\footnote{This lemma is
verified in Lean by the CSL core declaration listed in
\Cref{app:formalization}; the Lean theorem proves the four-clause
structure from the six structural premises (a)--(f) supplied as
parametric hypotheses.  The narrowing-disclosure discipline is recorded
in \Cref{app:formalization}.}
\end{proof}

This is the structural law on which the rest of the paper rests.
\Cref{sec:v-differential} classifies the four named substrates by the
kind of exposure their autonomous lifts can exhibit.
\Cref{sec:granularity-mismatch} then records why the typed quotients
used here do not project to a standard \(P\)/\(NP\) separation.  The
\CSL{} remains at \DiagnosisGrade{} throughout: it is the bounded
finite-carrier structural law, not an asymptotic theorem over all
polynomial-time histories.
 

\section{Substrate differentiation}
\label{sec:v-differential}

The \CSL{} applies to the framework-wide carrier described above, but
its consequences are substrate-differentiated.  Some substrates admit
autonomous lifts whose current observables expose witness-side content;
others admit autonomous lifts whose current observables expose only
trace state.  The V-Differential classification records this difference
across four named substrates.

An autonomous lift \(\Lift\) on a substrate \(\Substrate\) is
\emph{witness-content-exposing} when
\(\Exposure(\Lift,\Substrate)\) holds with a current observable that
depends on witness-side data: an assignment, a certificate value, or a
field that reveals a specific solution coordinate.  Equivalently, the
polynomial-time current observable reads through the lift state to a
witness-side coordinate.  This is a statement about the selected
carrier and its observation channel, not a claim that the corresponding
host problem has been solved.

By contrast, an autonomous lift is \emph{trace-state-exposing} when
\(\Exposure(\Lift,\Substrate)\) holds only through trace bookkeeping or
replay-side data: a step count, a partial decision-tree position, an
audit field, or a similar record of how the computation proceeded.
Such a lift is exposed in the sense of \Cref{sec:visibility-coupling},
because a current observable factors through the lift state.  It does
not thereby reveal witness-side content such as a satisfying assignment,
a regulator-prediction value, a weak-solution selection datum, or a
zeta-zero certificate coordinate.

\begin{proposition}[V-Differential classification]
\label{prop:v-differential}
Within the typed substrate apparatus, the following four named
substrates admit autonomous lifts of the indicated exposure type.
\begin{enumerate}
  \item \emph{BSD substrate.}  The Birch--Swinnerton-Dyer substrate
  \citep{BirchSwinnertonDyer1965NotesII}, read here through
  elliptic-curve regulator content
  \citep{Silverman2009ArithmeticEllipticCurves}, admits autonomous
  lifts that are witness-content-exposing.

  \item \emph{NS substrate.}  The Navier--Stokes substrate
  \citep{Fefferman2006NavierStokesClay}, read here through
  weak-solution selection content \citep{Leray1934LiquideVisqueux},
  admits autonomous lifts that are witness-content-exposing.

  \item \emph{PvNP substrate.}  The \(P\)-versus-\(NP\) substrate
  \citep{Cook2000PNPClay}, read here through satisfying-assignment
  witness content for \(3\)-SAT \citep{Cook1971TheoremProving}, admits
  autonomous lifts that are trace-state-exposing but not
  witness-content-exposing under the polynomial-time current instrument.

  \item \emph{RH substrate.}  The Riemann-hypothesis substrate
  \citep{Bombieri2000RiemannClay}, read here through zeta-zero
  certificate content \citep{Conrey2003RiemannHypothesis}, admits
  autonomous lifts that are trace-state-exposing but not
  witness-content-exposing under the polynomial-time current instrument.
\end{enumerate}
In each case the lift's existence and its classification are
framework-internal facts about the substrate's adequacy-residual
machinery.  The classification records which exposure channel,
witness-content or trace-state, is available under the framework's
lawful instrument.
\end{proposition}

\begin{proof}
The proof consumes, for each of the four named substrates, the
framework-internal exposure profile of its admissible autonomous lifts.
These are per-substrate inputs to the V-Differential classification,
supplied by the substrate's adequacy-residual machinery; they are not
derived inside this proof.

For the BSD substrate, the elliptic-curve regulator content carried by
lift internal-state readouts is, on the witness-content channel,
polynomial-time exposed to current observation by a deterministic
readout of the regulator coordinate.  Hence the BSD substrate admits
autonomous lifts that are witness-content-exposing.

For the NS substrate, the weak-solution selection content carried by
lift internal-state readouts admits, on the witness-content channel, a
polynomial-time current readout of the selection datum.  Hence the NS
substrate admits autonomous lifts that are witness-content-exposing.

For the PvNP substrate, the satisfying-assignment witness content for
\(3\)-SAT is, under the polynomial-time current instrument, available
only to predictive observation; current observation under that
instrument admits only the trace-state side, such as replay records,
decision-tree partial position, and audit fields.  Hence the PvNP
substrate admits autonomous lifts that are trace-state-exposing but not
witness-content-exposing under the polynomial-time current instrument.

For the RH substrate, the zeta-zero certificate content is, under the
polynomial-time current instrument, similarly available only to
predictive observation; current observation admits only trace state.
Hence the RH substrate admits autonomous lifts that are
trace-state-exposing but not witness-content-exposing under the
polynomial-time current instrument.

These four substrate placements record the kind of exposure their
admissible autonomous lifts admit.  The classification is structural:
it records what current observation can reveal under the framework's
lawful instrument, without making framework-internal claims about
\CSL{} dissolution or preservation for each substrate.\footnote{This proposition is
verified in Lean by the V-Differential declaration listed in
\Cref{app:formalization}; the Lean theorem proves the sorted
classification from per-substrate empirical hypotheses, namely
witness-exposure-fraction bounds for trace-state-exposing substrates and
witness-exposing lift constructions for witness-content-exposing
substrates, supplied as typed structural inputs.  The
narrowing-disclosure discipline is recorded in \Cref{app:formalization}.}
\end{proof}

This proposition is a classification of framework carriers, not a
theorem resolving the Birch--Swinnerton-Dyer conjecture, the
Navier--Stokes problem, the Riemann hypothesis, or the \(P\)-versus-\(NP\)
problem.  The per-substrate exposure profiles are framework-internal
inputs to the classification, supplied by the corresponding substrate's
adequacy-residual machinery; the classification does not derive the
host-mathematical substrate facts.

The V-Differential gives the exposure analysis substrate-level content.
For the PvNP substrate in particular, the trace-state-only
classification has structural consequences for what kind of
framework-internal conditional theorem could be formulated over the
\(P\)-versus-\(NP\) problem; those consequences fall outside the scope of
the present paper and are noted in \Cref{sec:scope,sec:discussion}.
The next section turns from substrate-level structure to the
relationship between the framework's typed quotients and the standard
classes \(P\) and \(NP\).
 

\section{Granularity mismatch}
\label{sec:granularity-mismatch}

The framework quotients \(\Qquotient\) and \(\Mquotient\) are quotients
of computation histories.  The standard complexity classes \(P\) and
\(NP\), by contrast, are language classes
\citep{AroraBarak2009Complexity}: they classify membership predicates on
strings, or verifier predicates whose existential projection is a
language \citep{Cook2000PNPClay}.  The natural maps from the typed quotients to
the standard language-level objects are therefore projections that
discard information carried by the history.

There are two projections to keep separate.  The map \(\PiQP\) sends a
current quotient class to the polynomial-time decision predicate read
from that class.  The map \(\PiMNP\) sends a predictive quotient class
to a verifier predicate on instance--witness pairs.  Both are natural
standard quotient constructions, and both are well-defined for the
same reason the quotient definitions were well-defined in
\Cref{sec:framework}: the relevant output is one of the observables
used to form the quotient.  Neither projection is a bijection.  The
two obstruction propositions below record the mismatch; the next
subsection gives the constructive witness showing that strict
\(\Mquotient\)-over-\(\Qquotient\) refinement can collapse to one
language-level outcome under these projections.

\begin{definition}[Q-to-P projection]
\label{def:projection-q-to-p}
The \emph{Q-to-P projection} is the map
\[
  \PiQP:\Qquotient\longrightarrow
  \{\text{polynomial-time decision predicates}\}
\]
that sends a current class \(\Qclass{h}\) to the polynomial-time
decision predicate read off from any representative history \(h\) in
that class.  This is well-defined because the decision output of a
history is itself a polynomial-time current observable
\(\CurrentObservable\), hence is constant on \(\Qquotient\)-classes.
\end{definition}

\begin{proposition}[\(\PiQP\) is not a bijection onto \(P\)]
\label{prop:c1-not-bijective}
The projection \(\PiQP\) is not a bijection from \(\Qquotient\) onto
the standard class \(P\).  The obstruction has three parts.
\begin{enumerate}
  \item \emph{Histories versus languages.}  A \(\Qquotient\)-class is
  a class of histories, possibly over a single bounded instance.  A
  polynomial-time decision predicate is a function on all inputs.  A
  single current class does not by itself define a language predicate
  over all input strings.

  \item \emph{Trace observables beyond decision output.}  The quotient
  \(\Qquotient\) identifies histories by equality of all declared
  polynomial-time current observables, not merely by equality of the
  decision bit.  Two histories computing the same language may still be
  \(\Qquotient\)-distinct because their trace, solver-prefix, or audit
  bookkeeping observables differ.

  \item \emph{Machine multiplicity.}  Different polynomial-time
  machines may decide the same language while producing different
  \(\Qquotient\)-observable histories on the same instance.  These
  histories can map to different current classes but to the same
  language predicate.
\end{enumerate}
\end{proposition}

\begin{proof}
Each item is a direct comparison of the two granularities.  Item (1)
shows that a history class is not, without extra data over all inputs,
a language predicate.  Items (2) and (3) show non-injectivity after
projection to language-level predicates: \(\Qquotient\) retains trace,
solver-prefix, audit, and machine-choice information that standard
language equality forgets.  Thus a theorem-grade bijection with
standard \(P\) is unavailable without an additional quotient that
forgets those features.  Equivalently, once any of the carrier-side
sources above supplies two distinct \(\Qquotient\)-classes with the same
\(\PiQP\)-image, \(\PiQP\) fails to be injective, and hence cannot be a
bijection onto \(P\).\footnote{This proposition is verified in Lean by
the C1-obstruction collision-schema declaration listed in
\Cref{app:formalization}; the Lean theorem proves a collision schema
parameterized by a supplied collision witness, leaving the carrier-side
obstruction clauses as paper-side input.  The narrowing-disclosure
discipline is recorded in \Cref{app:formalization}.}
\end{proof}

\begin{definition}[M-to-NP projection]
\label{def:projection-m-to-np}
The \emph{M-to-NP projection} is the map
\[
  \PiMNP:\Mquotient\longrightarrow
  \{\text{verifier predicates }V(x,w)\text{ on instance--witness pairs}\}
\]
that sends a predictive class \(\Mclass{h}\) to the verifier predicate
observed in that class.  This is well-defined because verifier
acceptance is a witness-augmented polynomial-time observable, hence is
constant on \(\Mquotient\)-classes.
\end{definition}

\begin{proposition}[\(\PiMNP\) is not a bijection onto \(NP\)]
\label{prop:c2-not-bijective}
The projection \(\PiMNP\) is not a bijection from \(\Mquotient\) onto
the standard class \(NP\).  The obstruction is witness identity:
\(\Mquotient\) can distinguish different satisfying witnesses for the
same instance, through witness-trace features such as Hamming-distance
fingerprints, satisfaction-fraction sequences, and partial-assignment
branching records \citep{BiereHeuleMaarenWalsh2009HandbookSAT}.  Standard \(NP\)-membership is existential; it
retains only that some witness exists, not which witness or which
witness trace was observed.
\end{proposition}

\begin{proof}
Consider any substrate that admits a satisfiable instance with two
distinct witnesses, as in the usual SAT example with multiple
satisfying assignments.  The two corresponding histories can be
separated by witness-augmented observables, so they can represent
different \(\Mquotient\)-classes.  Under standard \(NP\)-membership,
however, both witnesses certify the same existential fact for the same
instance.  The projection therefore forgets the witness identity that
\(\Mquotient\) records, and a bijection onto standard \(NP\) cannot be
maintained without quotienting away that identity.  Equivalently, once
such a witness-identity collision supplies two distinct
\(\Mquotient\)-classes with the same \(\PiMNP\)-image, \(\PiMNP\) fails
to be injective, and hence cannot be a bijection onto \(NP\).\footnote{This
proposition is verified in Lean by the C2-obstruction collision-schema
declaration listed in \Cref{app:formalization}; the Lean theorem proves
a collision schema parameterized by a supplied collision witness, with
the witness-identity content paper-side.  The narrowing-disclosure
discipline is recorded in \Cref{app:formalization}.}
\end{proof}

Thus \(\PiQP\) and \(\PiMNP\) become surjective onto natural
language-level target classes only after additional quotienting that
forgets trace and witness identity.  The full granularity-mismatch
theorem in the next subsection records the consequence: strict
\(\Mquotient\)-over-\(\Qquotient\) refinement at the history level can
collapse under projection to a single language-level outcome.  That
theorem is stated at \RDiagnosis{}.

With the projections and non-bijection records in place, the
granularity-mismatch theorem can be stated precisely.  It is an
existence theorem over the typed carrier: the witness is a pair of
histories whose predictive classes differ while their current class and
language-level projections agree.

\begin{theorem}[Granularity mismatch]
\label{thm:granularity-mismatch}
There exist a substrate \(\Substrate\), an autonomous lift \(\Lift\) on
\(\Substrate\), and histories \(h,h'\in\HistorySpace\) such that
\begin{enumerate}
  \item \(\Mclass{h}\neq\Mclass{h'}\);
  \item \(\Qclass{h}=\Qclass{h'}\);
  \item \(\PiMNP(\Mclass{h})=\PiMNP(\Mclass{h'})\);
  \item \(\PiQP(\Qclass{h})=\PiQP(\Qclass{h'})\).
\end{enumerate}
Thus strict \(\Mquotient\)-over-\(\Qquotient\) refinement at the
history level can fail to survive projection to standard
language-level classes.  The strictness lives at the witness-trace
level, which the projections discard.
\end{theorem}

\begin{proof}
We construct the required data.  Let the substrate \(\Substrate\) be a
SAT-like bounded substrate containing a single \(3\)-SAT instance
\(\phi_{0}\) with two distinct satisfying assignments.  Let the
autonomous lift \(\Lift\) be any feedback-coupled or
planted-suppressed lift whose histories record both the fixed instance
and the satisfying assignment used by the lift.  On the finite carrier
under consideration, model histories as pairs
\[
  (\text{instance},\text{witness bit})
  \qquad\text{with instance fixed at }\phi_{0}
  \text{ and witness bit in }\{0,1\}.
\]
Let \(h=(\phi_{0},0)\) and \(h'=(\phi_{0},1)\).  We verify the four
asserted conditions.

First, \(\Mclass{h}\neq\Mclass{h'}\).  The two histories differ in the
witness bit, and witness-augmented observation includes observables
that may read witness-side data.  The observable that reads this bit
separates \(h\) from \(h'\), so the two histories lie in distinct
\(\Mquotient\)-classes.

Second, \(\Qclass{h}=\Qclass{h'}\).  The current observables in this
carrier record the fixed instance, the decision-level output, and
trace data that do not read the witness bit.  Since \(h\) and \(h'\)
share the same instance \(\phi_{0}\) and have the same decision output,
namely that \(\phi_{0}\) is satisfiable, every declared current
observable takes the same value on the two histories.  Hence they are
\(\Qquotient\)-equivalent.

Third, \(\PiMNP(\Mclass{h})=\PiMNP(\Mclass{h'})\).  Although the two
\(\Mquotient\)-classes remember which satisfying assignment was used,
the \(M\)-to-\(NP\) projection forgets witness identity and retains the
existential verifier-level fact.  Both witnesses certify the same
statement, namely that the fixed instance \(\phi_{0}\) is satisfiable.
Thus the two predictive classes have the same image under \(\PiMNP\).

Fourth, \(\PiQP(\Qclass{h})=\PiQP(\Qclass{h'})\).  This is immediate
from the equality of current classes, and concretely both histories
project to the same polynomial-time decision predicate on the fixed
instance: the predicate answering yes for the satisfiable instance
\(\phi_{0}\).

The constructed \(\Substrate\), \(\Lift\), \(h\), and \(h'\) therefore
satisfy all four clauses in the theorem statement.  The strict
refinement is real inside the typed history carrier, because the
predictive quotient separates the witness-bit histories, but it is
invisible to the standard language-level projections, because both
projections discard that witness identity.\footnote{This theorem is
verified in Lean by the granularity-mismatch declaration listed in
\Cref{app:formalization}; the Lean theorem exhibits a SAT-like
finite-carrier witness with \(\Mquotient\) instantiated as literal
history equality, rather than as the kernel of a witness-augmented
observable.  The narrowing-disclosure discipline is recorded in
\Cref{app:formalization}.}
\end{proof}

The theorem is at \RDiagnosis{}.  In plain terms, the structural
obstruction is identified and supported by a constructive finite-carrier
witness, while the bridging representation theorem from arbitrary
polynomial-time deciders to autonomous-lift carriers preserving
witness-content exposure is named but not constructed here.  The
theorem records the obstruction; it does not claim that all
polynomial-time-decider settings exhibit the obstruction asymptotically.

The granularity-mismatch theorem completes the substantive body of the
paper.  The next section states the nonclaims explicitly, including
what this paper does not prove about standard \(P\) and \(NP\);
see \Cref{sec:scope}.
 

\section{Scope and non-claims}
\label{sec:scope}

The body has stated the \CSL{} at \DiagnosisGrade{}, the
V-Differential substrate classification, and the granularity-mismatch
theorem at \RDiagnosis{}.  This section gives an explicit accounting
of what those results do and do not claim, so that the preceding
statements are not mistaken for standard complexity-theoretic
conclusions.

\begin{remark}[Nonclaim: no standard \(P\neq NP\) consequence]
The \CSL{} is a framework-internal structural law at diagnosis grade
over predictive-quotient holonomy on autonomous-lifted polynomial-time
substrates; it does not imply standard \(P\neq NP\).  This paper does
not claim, and the present results do not provide, a standard-ZFC
unconditional proof of \(P\neq NP\)
\citep{Cook2000PNPClay,Fortnow2009StatusPNP}.  The granularity-mismatch
theorem records the precise reason that the framework-internal strict
refinement does not project to a standard \(P\)/\(NP\) separation.
\end{remark}

\begin{remark}[Nonclaim: no asymptotic theorem-grade CSL]
The \CSL{} is hardened to diagnosis grade by the cascade work it
summarizes: bounded finite-carrier evidence with falsifier audit.  It
is not stated as an asymptotic theorem over all polynomial-time
computation histories.  The empirical hardening is a paper-side
certification rather than a derivable consequence of the lemma's formal
content; the boundary is set explicitly by the \DiagnosisGrade{}
qualifier.
\end{remark}

\begin{remark}[Nonclaim: no representation theorem]
The granularity-mismatch theorem at R-Diagnosis identifies the
structural obstruction; the representation theorem is named but not
constructed.  More explicitly, the theorem identifies the structural
obstruction between the framework's typed quotients and the standard
\(P\)/\(NP\) language classes, and it exhibits a constructive
finite-carrier witness for the existential conclusion.  The bridging
representation theorem from arbitrary polynomial-time deciders to
autonomous-lift carriers preserving witness-content exposure is the
asymptotic step this paper does not establish; the \RDiagnosis{} grade
label records this gap explicitly.
\end{remark}

\begin{remark}[Additional nonclaims]
The framework is similar in spirit to, but structurally distinct from,
the natural-proofs barrier of Razborov--Rudich
\citep{RazborovRudich1997NaturalProofs}, the algebrization barrier of
Aaronson--Wigderson \citep{AaronsonWigderson2009Algebrization}, and the
relativization barrier of Baker--Gill--Solovay
\citep{BakerGillSolovay1975Relativization}.  The predictive-quotient
holonomy carrier is not a circumvention of any of those three barriers
in their native senses; the relationship is discussed in
\Cref{sec:discussion}.

The framework also uses standard mathematical constructions throughout:
quotients of finite carriers, projection maps, the equivalence of
refinement and surjection between quotient relations, and finite
Boolean carriers for the granularity-mismatch witness.  The
contribution is the typed \CSL{} apparatus and the
granularity-mismatch packaging, not these underlying mathematical
constructions.
\end{remark}

Positively, the paper claims three things.  First, it states a typed
\CSL{} over the predictive-quotient holonomy carrier on
autonomous-lifted polynomial-time substrates, with dissolution,
preservation, degenerate, and strength-independence clauses, at
\DiagnosisGrade{}.  Second, it gives the V-Differential classification
of four named substrates by whether their autonomous lifts expose
witness content or only trace state.  Third, it proves the
granularity-mismatch theorem at \RDiagnosis{}, with a constructive
finite-carrier witness showing how history-level strictness can
collapse under standard language-level projection.

The next section discusses how these results relate to classical
complexity-theoretic barriers and to neighboring framework directions.
 

\section{Discussion}
\label{sec:discussion}

The preceding sections give a framework-internal structural law at
\DiagnosisGrade{} and a projection obstruction.  This discussion
positions those results against classical complexity barriers, comments
on the minimal finite witness, the standard constructions used in the
paper, and the word ``candidate'' in \CSLcore{}, and points to
neighboring framework directions.

\paragraph{Classical complexity barriers.}
The framework is similar in spirit to, but structurally distinct from,
the natural-proofs barrier of Razborov--Rudich
\citep{RazborovRudich1997NaturalProofs}, the algebrization barrier of
Aaronson--Wigderson \citep{AaronsonWigderson2009Algebrization}, and the
relativization barrier of Baker--Gill--Solovay
\citep{BakerGillSolovay1975Relativization}; each constrains a different
class of proof techniques inside its own setting
\citep{Wigderson2019MathematicsComputation,Aaronson2016PNP}.  The
\CSL{}'s load-bearing variable is the binary exposure predicate
\(\Exposure\), and its strength-independence clause is explicit,
whereas the natural-proofs barrier is about cryptographic strength
\citep{GoldreichGoldwasserMicali1986RandomFunctions,HastadImpagliazzoLevinLuby1999PRG}.
The predictive-quotient holonomy carrier does not invoke oracle access:
quotient classes depend on the declared families of polynomial-time
current and witness-augmented observables, so a relativized or
algebrized version would change the observable families and therefore
change the quotients themselves.  The framework is therefore neither a
circumvention nor a refinement of these three barriers; it is a
carrier-internal statement with a different dependency structure.

\paragraph{The minimal witness.}
The constructive finite-carrier witness in
\Cref{thm:granularity-mismatch} is intentionally minimal: one bounded
instance and one Boolean witness bit.  This is the smallest carrier on
which the structural-obstruction shape is visible: \(\Mquotient\)
strictly refines \(\Qquotient\) because of witness identity, while the
language-level projections forget that identity.  A non-minimal
witness, for example SAT under a feedback-coupled autonomous lift with
several satisfying assignments, carries the same structural shape.
The \RDiagnosis{} label records this scope: the theorem captures the
typed obstruction and does not promote the finite witness into a deep
asymptotic complexity claim.

\paragraph{Standard constructions.}
The paper uses standard mathematical constructions throughout:
quotients \citep{Halmos1960NaiveSetTheory},
projection maps, refinement-as-surjection comparisons, and finite
Boolean carriers.  The contribution is the typed \CSL{}
apparatus and the granularity-mismatch packaging, not novelty in those
underlying constructions.  The audited finite-carrier framework of
Foundations II and the finite-audit interaction calculus of
Foundations III provide the typed setting in which the quotients are
audited, while the predictive-quotient holonomy carrier was developed
in the holonomy-with-memory work
\citep{TsiokosFoundationsII2026,TsiokosFoundationsIII2026,TsiokosHolonomyMemory2026}.
The \CSL{} is a structural law at \DiagnosisGrade{} on that carrier.

\paragraph{Why ``candidate''.}
The law is labeled \CSLcore{} rather than ``Structural Theorem''
because it is hardened to diagnosis grade by the cascade work it
summarizes: bounded finite-carrier evidence with a falsifier audit, not
an asymptotic theorem over all polynomial-time computation histories.
The word ``candidate'' marks this epistemic placement.  The
four-clause structural shape is established by typed reasoning from the
structural premises stated in \Cref{sec:csl-core}; the universal claim
over all asymptotic carriers is a separate theorem-grade target.  The
label marks the work needed before the law could be promoted.

\paragraph{Open directions.}
The V-Differential classification places PvNP and RH as trace-state
exposing, while BSD and NS are witness-content exposing.  These are
framework-internal structural placements, not claims about resolving
the \(P\)-versus-\(NP\) problem, the Riemann hypothesis, the
Birch--Swinnerton-Dyer conjecture, or the Navier--Stokes problem.  They
suggest research directions: one may ask what additional
substrate-specific typed content would be needed to turn each
classification into a conditional framework theorem.  The present
paper does not develop those directions.

The conclusion summarizes the contribution and the extension targets
left by this scope.
 

\section{Conclusion}
\label{sec:conclusion}

This paper has stated a typed \CSLcore{} (\(\CSL\)) over the
predictive-quotient holonomy carrier on autonomous-lifted
polynomial-time substrates.  It has also proved a
granularity-mismatch theorem recording why the framework's typed
quotients are strictly finer than the standard classes \(P\) and
\(NP\) at the history level: strictness can live in witness-trace
identity and disappear under standard language-level projection.  The
\CSL{} remains at \DiagnosisGrade{}, and the granularity-mismatch
theorem remains at \RDiagnosis{}.  Together with the V-Differential
classification of four named host substrates, the headline is a typed
structural picture, not a standard-ZFC theorem about \(P\)-versus-\(NP\).

Three kinds of work remain.  First, asymptotic hardening of the
\CSL{} would require moving beyond bounded finite carriers to a
theorem-grade statement over polynomial-time computation histories
\citep{AroraBarak2009Complexity,Wigderson2019MathematicsComputation}.
Second, the \RDiagnosis{} grade of the granularity-mismatch theorem
marks the missing representation theorem: a construction from
arbitrary polynomial-time deciders to autonomous-lift carriers that
preserve witness-content exposure.  Third, the V-Differential
classification suggests substrate-specific conditional frameworks: for
each named substrate, one can ask what additional typed content would
turn its trace-state-only or witness-content-exposing placement into a
conditional framework theorem.  These are targets named by the present
paper, not results established here.

The present paper therefore does not claim a standard-ZFC
unconditional proof of \(P\neq NP\).  It does not claim circumvention
of the natural-proofs, algebrization, or relativization barriers in
their native senses.  It also does not claim resolutions of the host
problems named by the V-Differential classification.  The grade
qualifiers \(\DiagnosisGrade{}\) and \(\RDiagnosis{}\) record the
binding scope of the two central claims.

The framework provides a typed structural vocabulary for naming when
prediction is more than current observation, when it collapses to
current observation, and why that distinction does not project to the
standard complexity-theoretic story without further work.
 
\appendix


\section{Formalization (Lean coverage)}
\label{app:formalization}

This appendix records how the paper-side definitions and claims are
realized in Lean \citep{DeMouraUllrich2021Lean4}.  It is a per-row map
from the labeled mathematical statements in the body to their
declarations in the \texttt{SixBirdsHiddenness.Hiddenness} namespace,
together with the coverage and alignment language used in the paper.  It is not a
restatement of the body theorems, not a substitute for the mathematical
proof sketches, and not an external-bibliography list.

Lean identifier strings appear here and only here.  The body prose uses
paper-prose names such as ``current quotient'' or ``granularity
mismatch''; this appendix gives the corresponding module path and
declaration name for audit and review.

\paragraph{Lean-disclosure wording.}
The body and appendix use the following wording discipline.
\begin{center}
\small
\begin{tabularx}{\textwidth}{@{}>{\raggedright\arraybackslash}p{0.22\textwidth}>{\raggedright\arraybackslash}p{0.24\textwidth}>{\raggedright\arraybackslash}X@{}}
\toprule
Coverage / alignment & Permitted wording & Meaning in this paper \\
\midrule
theorem / faithful & ``Lean proves'' or ``verified in Lean'' & The Lean theorem matches the paper statement. \\
theorem / narrowed surrogate & ``Lean proves the stated scope; the broader claim is paper-side with narrowing disclosed here'' & The structural shape is mechanized, with explicit premises or finite witnesses exposed. \\
definition / faithful & ``realized in Lean as \texttt{Module.name}'' & The Lean declaration is the paper definition's typed realization. \\
definition / narrowed surrogate & ``realized in Lean with the parameter or scope exposed'' & The Lean definition exposes an extra structural field or predicate. \\
not mechanized & No Lean citation & The row is a nonclaim; the body uses the scope-fence wording. \\
\bottomrule
\end{tabularx}
\end{center}

\begin{landscape}
\footnotesize
\setlength{\tabcolsep}{2pt}
\renewcommand{\arraystretch}{0.88}
\noindent\textbf{Coverage table.}\par\smallskip
\begin{longtable}{@{}>{\raggedright\arraybackslash}p{0.19\linewidth}>{\raggedright\arraybackslash}p{0.12\linewidth}>{\raggedright\arraybackslash}p{0.36\linewidth}>{\raggedright\arraybackslash}p{0.09\linewidth}>{\raggedright\arraybackslash}p{0.21\linewidth}@{}}
\toprule
Paper label & Paper-prose name & Lean decl & Coverage / alignment & Disclosure wording (body location) \\
\midrule
\endfirsthead
\toprule
Paper label & Paper-prose name & Lean decl & Coverage / alignment & Disclosure wording (body location) \\
\midrule
\endhead
\midrule
\multicolumn{5}{r}{\footnotesize\itshape continued on next page} \\
\endfoot
\bottomrule
\endlastfoot
\path|def:hiddenness:substrate| & substrate & \mbox{\texttt{Hiddenness.HistorySpace.substrate}} & definition / faithful & Definition entry; realized in Lean as listed here (\Cref{sec:framework}). \\
\path|def:hiddenness:history-space| & history space & \mbox{\texttt{Hiddenness.HistorySpace.historySpace}} & definition / faithful & Definition entry; realized in Lean as listed here (\Cref{sec:framework}). \\
\path|def:hiddenness:current-quotient| & current quotient & \mbox{\texttt{Hiddenness.Quotients.currentQuotient}} & definition / faithful & Definition entry; realized in Lean as listed here (\Cref{sec:framework}). \\
\path|def:hiddenness:predictive-quotient| & predictive quotient & \mbox{\texttt{Hiddenness.Quotients.predictiveQuotient}} & definition / faithful & Definition entry; realized in Lean as listed here (\Cref{sec:framework}). \\
\path|def:hiddenness:forget-witness-map| & forget-witness map & \mbox{\texttt{Hiddenness.Quotients.forgetWitnessMap}} & definition / faithful & Definition entry; realized in Lean as listed here (\Cref{sec:framework}). \\
\path|prop:hiddenness:witness-refinement| & witness refinement proposition & \mbox{\texttt{Hiddenness.Quotients.witnessRefinement}} & theorem / faithful & Verified in Lean; body proof explains the current-observable inclusion (\Cref{sec:framework}). \\
\path|def:hiddenness:autonomous-lift| & autonomous lift & \mbox{\texttt{Hiddenness.AutonomousLift.autonomousLift}} & definition / narrowed surrogate & Realized with family-characterizing predicates exposed in the Lean record (\Cref{sec:autonomous-lifts}). \\
\path|def:hiddenness:lift-family-A| & Lift A (self-generating substrate) & \mbox{\texttt{Hiddenness.AutonomousLift.liftFamilyA}} & definition / faithful & Definition entry; realized in Lean as listed here (\Cref{sec:autonomous-lifts}). \\
\path|def:hiddenness:lift-family-B| & Lift B (feedback-coupled) & \mbox{\texttt{Hiddenness.AutonomousLift.liftFamilyB}} & definition / faithful & Definition entry; realized in Lean as listed here (\Cref{sec:autonomous-lifts}). \\
\path|def:hiddenness:lift-family-C| & Lift C (internally-seeded random) & \mbox{\texttt{Hiddenness.AutonomousLift.liftFamilyC}} & definition / faithful & Definition entry; realized in Lean as listed here (\Cref{sec:autonomous-lifts}). \\
\path|def:hiddenness:lift-family-D| & Lift D (planted-suppressed internalization) & \mbox{\texttt{Hiddenness.AutonomousLift.liftFamilyD}} & definition / faithful & Definition entry; realized in Lean as listed here (\Cref{sec:autonomous-lifts}). \\
\texttt{def:hiddenness:\allowbreak trace-fingerprint-\allowbreak loop-class} & trace-fingerprint loop class & \mbox{\texttt{Hiddenness.LoopAction.traceFingerprintLoopClass}} & definition / faithful & Definition entry; realized in Lean as listed here (\Cref{sec:loop-action}). \\
\path|def:hiddenness:loop-action-q| & loop action on Q & \mbox{\texttt{Hiddenness.LoopAction.loopActionQ}} & definition / faithful & Definition entry; realized in Lean as listed here (\Cref{sec:loop-action}). \\
\path|def:hiddenness:loop-action-m| & loop action on M & \mbox{\texttt{Hiddenness.LoopAction.loopActionM}} & definition / faithful & Definition entry; realized in Lean as listed here (\Cref{sec:loop-action}). \\
\path|def:hiddenness:asymmetry| & M-vs-Q loop-action asymmetry & \mbox{\texttt{Hiddenness.LoopAction.asymmetry}} & definition / faithful & Definition entry; realized in Lean as listed here (\Cref{sec:loop-action}). \\
\path|def:hiddenness:visibility-coupling| & visibility coupling / Exposure predicate & \mbox{\texttt{Hiddenness.VisibilityCoupling.visibilityCoupling}} & definition / faithful & Definition entry; realized with realized-support nonconstancy (\Cref{sec:visibility-coupling}). \\
\path|def:hiddenness:predictive-witness-content| & predictive witness content & \mbox{\texttt{Hiddenness.WitnessContent.predictiveWitnessContent}} & definition / faithful & Definition entry; realized in Lean as listed here (\Cref{sec:witness-content}). \\
\path|lem:csl-core| & CSL core lemma / Candidate Structural Law & \mbox{\texttt{Hiddenness.CSLCore.cslCore}} & theorem / narrowed surrogate & Lean proves the four clauses from six structural premises; empirical hardening is paper-side (\Cref{sec:csl-core}). \\
\path|prop:hiddenness:v-differential| & V-Differential proposition & \mbox{\texttt{Hiddenness.VDifferential.vDifferential}} & theorem / narrowed surrogate & Lean proves the sorted classification from per-substrate hypotheses (\Cref{sec:v-differential}). \\
\path|def:hiddenness:projection-q-to-p| & Q-to-P projection & \mbox{\texttt{Hiddenness.GranularityMismatch.projectionQToP}} & definition / faithful & Definition entry; realized in Lean as listed here (\Cref{sec:granularity-mismatch}). \\
\path|prop:hiddenness:c1-not-bijective| & C1 obstruction (Q-to-P non-bijection) & \mbox{\texttt{Hiddenness.GranularityMismatch.c1NotBijective}} & theorem / narrowed surrogate & Lean proves a collision schema: any supplied collision between distinct \(\Qquotient\)-classes sharing a \(\PiQP\) image yields non-injectivity. The three carrier-side obstruction clauses are paper-side (\Cref{sec:granularity-mismatch}). \\
\path|def:hiddenness:projection-m-to-np| & M-to-NP projection & \mbox{\texttt{Hiddenness.GranularityMismatch.projectionMToNP}} & definition / faithful & Definition entry; realized in Lean as listed here (\Cref{sec:granularity-mismatch}). \\
\path|prop:hiddenness:c2-not-bijective| & C2 obstruction (M-to-NP non-bijection) & \mbox{\texttt{Hiddenness.GranularityMismatch.c2NotBijective}} & theorem / narrowed surrogate & Lean proves a collision schema: any supplied collision between distinct \(\Mquotient\)-classes sharing a \(\PiMNP\) image yields non-injectivity. The witness-identity content is paper-side (\Cref{sec:granularity-mismatch}). \\
\path|thm:granularity-mismatch| & granularity-mismatch theorem & \mbox{\texttt{Hiddenness.GranularityMismatch.granularityMismatch}} & theorem / narrowed surrogate & Lean exhibits a SAT-like finite-carrier witness with the witness-discriminating setoid \(M\) instantiated as literal history equality and a constant verifier output, rather than as the kernel of a witness-augmented observable reading the witness bit; no standard \(P\)-vs-\(NP\) bridge is claimed (\Cref{sec:granularity-mismatch}). \\
\path|thm:standard-pvnp-implication| & standard \(P\)-vs-\(NP\) implication nonclaim & --- & not mechanized & No Lean citation; scope-fence wording only (\Cref{sec:scope}). \\
\path|thm:csl-theorem-grade-asymptotic| & asymptotic theorem-grade CSL nonclaim & --- & not mechanized & No Lean citation; diagnosis-grade scope wording only (\Cref{sec:scope}). \\
\path|thm:razborov-rudich-internal-restatement| & classical-barrier circumvention nonclaim & --- & not mechanized & No Lean citation; barrier-comparison wording only (\Cref{sec:scope,sec:discussion}). \\
\end{longtable}
\end{landscape}

The three nonclaim rows carry no Lean coverage.  They are included in
the table because they are paper-side rows of record, but the body does
not cite Lean for them.  Their content is expressed through the
scope-fence language in \Cref{sec:scope}: no standard-ZFC unconditional
\(P\neq NP\) consequence, no asymptotic theorem-grade \CSL{}, and no
native circumvention claim for the classical barriers discussed in
\Cref{sec:discussion}.

The narrowed-surrogate rows are the main disclosure points.  For
\path|lem:csl-core|, Lean proves the four-clause structure from the six
structural premises stated in \Cref{sec:csl-core}; the
\DiagnosisGrade{} empirical hardening is paper-side.  For
\path|prop:hiddenness:v-differential|, Lean proves the sorted
classification from per-substrate empirical hypotheses.  For
\path|prop:hiddenness:c1-not-bijective| and
\path|prop:hiddenness:c2-not-bijective|, Lean proves a collision
schema: any supplied collision between distinct classes that share a
projection image yields non-injectivity.  The carrier-side obstruction
clauses (Q-to-P: histories-vs-languages, trace observables, machine
multiplicity; M-to-NP: witness-identity content) are paper-side; the
schemata apply once such a collision is supplied.  For
\path|thm:granularity-mismatch|, Lean realizes the typed projections
and a SAT-like finite-carrier witness; it does not claim a bridge from
that framework-internal projection to a standard \(P\)-vs-\(NP\)
implication.

The Lean development is mathlib-free
\citep{MathlibCommunity2020LeanMathematicalLibrary}.  All declarations
listed here live in the Six Birds Hiddenness hierarchy, with the
\path|Hiddenness| subnamespace reserved for this paper.  The theorem
rows pass the project's axiom audit under the trust base
\(\{\)\path|Classical.choice|, \path|propext|, \path|Quot.sound|,
\path|choice|, \path|funext|\(\}\) \citep{MartinLof1984IntuitionisticTypeTheory}.
 
\clearpage
\bibliographystyle{plainnat}
\bibliography{references}

\end{document}
